\documentclass{article}
\usepackage[UTF8]{ctex}
\usepackage{amsmath,amssymb,amsfonts}
\usepackage{enumitem}
\usepackage{geometry}
\geometry{a4paper, margin=1in}
% === 错题本排版优化（不参与题目标识） ===
\setlength{\parindent}{0pt}
\setlength{\parskip}{0.35em plus 0.10em minus 0.08em}
\setlength{\emergencystretch}{3em}
\setlist{itemsep=0.25em, topsep=0.40em, parsep=0pt}
\allowdisplaybreaks[2]
% === 错题本排版优化结束 ===
\begin{document}
\section*{西南财经大学2025年高等代数真题}
\section*{165．西南财经大学2025年高等代数试卷}

\subsection{题目1}
1．求 \(f(x)=2 x^{4}+x^{3}-2 x-1\) 的有理数解和在复数域上的因式分解．

\subsection{题目2}
2．线性方程组

\[
\left\{\begin{array}{l}
x_{1}+x_{2}+x_{3}+x_{4}=1 \\
x_{1}+2 x_{2}+3 x_{4}=2 \\
x_{1}+2 x_{2}+a x_{3}+3 x_{4}=b \\
x_{1}+3 x_{2}-x_{3}+(a+5) x_{4}=3
\end{array}\right.
\]

当 \(a, b\) 取何值时，该线性方程组有唯一解？无解？有无穷解？有无穷解的通解是什么？

\subsection{题目3}
3．向量组（I）：\(\alpha_{1}, \alpha_{2}, \cdots, \alpha_{m}\) ，（II）：\(\alpha_{1}, \alpha_{2}, \cdots, \alpha_{m}, \alpha_{m+1}, \cdots, \alpha_{n}\) 。证明：向量组（I）与（II）的秩相同的充分必要条件是向量组（II）可由（I）线性表出．

\subsection{题目4}
4．\(A=\left(\begin{array}{rrr}-1 & 1 & 0 \\ -2 & 2 & 0 \\ a & 3 & b\end{array}\right), \alpha_{1}, \alpha_{2}, \alpha_{3}\) 线性无关，且

\[
A \alpha_{1}=\alpha_{1}, A \alpha_{2}=\alpha_{2}, A \alpha_{3}=2 \alpha_{2} .
\]

证明：\(A\) 与 \(B=\left(\begin{array}{lll}1 & 0 & 0 \\ 0 & 1 & 2 \\ 0 & 0 & 0\end{array}\right)\) 相似，并计算出 \(a, b\) ．

\subsection{题目5}
5．已知

\[
A_{1}=\left(\begin{array}{cc}
2 & 1 \\
-1 & 3
\end{array}\right), A_{2}=\left(\begin{array}{ll}
1 & 0 \\
2 & 0
\end{array}\right), A_{3}=\left(\begin{array}{ll}
3 & 1 \\
1 & 3
\end{array}\right), A_{4}=\left(\begin{array}{cc}
1 & 1 \\
-3 & 3
\end{array}\right)
\]

求由其生成的子空间的一组基以及维数。

\subsection{题目6}
6．\(A\) 是 \(n\) 阶实对称正定矩阵，\(\alpha_{1}, \alpha_{2}, \cdots, \alpha_{n}, \beta\) 是 \(n\) 维欧氏空间 \(\mathbb{R}^{n}\) 中的 \(n+1\) 个向量。满足：

\begin{itemize}
  \item[（1）] \(\alpha_{i} \neq 0, i=1,2, \cdots, n\) ．
  \item[（2）] \(\alpha_{i}^{\mathrm{T}} A \alpha_{j}=0, i, j=1,2, \cdots, n, i \neq j\) ．
  \item[（3）] \(\beta\) 与 \(\alpha_{j}\) 正交，\(j=1,2, \cdots, n\) ．
\end{itemize}

证明：\(\beta=0\) ．

\subsection{题目7}
7．\(\sigma, \tau\) 是有限维线性空间中的线性变换．判断

\begin{itemize}
  \item[（1）] \(\operatorname{dim} \operatorname{Ker}(\sigma \tau) \leq \operatorname{dim} \operatorname{Ker} \sigma+\operatorname{dim} \operatorname{Ker} \tau\) ．
  \item[（2）] \(\operatorname{dim} \operatorname{Ker}(\sigma \tau)=\operatorname{dim} \operatorname{Ker}(\tau \sigma)\) ．
\end{itemize}

是否正确？并证明．

\subsection{题目8}
8．已知 \(n\) 阶实矩阵

\[
A=\left(\begin{array}{llll}
\lambda & 1 & & \\
& \lambda & 1 & \\
& & \lambda & \\
& & & 1 \\
& & & \lambda
\end{array}\right)
\]

记 \(W=\{B \mid A B=B A, B\) 为 \(n\) 阶实矩阵 \(\}\) 。证明：\(W\) 是线性空间，以及其维数为 \(n\) 。
\end{document}
